\section{Empirical Results}
\label{sec:empirical}

\subsection{Multi-State Comparison}

Table~\ref{tab:decl-baseline} reports Declination for \textsc{bisect} and enacted
2022 congressional maps across four states using VEST 2020 presidential returns.

\begin{table}[ht]
\centering
\caption{Declination: Bisect vs.\ Enacted 2022 Congressional Maps}
\label{tab:decl-baseline}
\begin{threeparttable}
\begin{tabular}{@{}llrrr@{}}
\toprule
State ($k$) & $\delta$ (Bisect$^\dagger$) & $\delta$ (Enacted) & $|\Delta\delta|$ & Reduction \\
\midrule
NC ($k=14$) & $+0.08$ & $+0.35$ & $0.27$ & 77\% \\
WI ($k=8$)  & $+0.07$ & $+0.28$ & $0.21$ & 75\% \\
TX ($k=38$) & $+0.09$ & $+0.40$ & $0.31$ & 78\% \\
FL ($k=28$) & $+0.08$ & $+0.32$ & $0.24$ & 75\% \\
\bottomrule
\end{tabular}
\begin{tablenotes}
\small
\item $^\dagger$ Single-run \textsc{bisect} result. Positive $\delta$ indicates
Republican advantage through packing/cracking.
\item Source: VEST 2020 presidential returns; 2020 Census tract geometry.
\end{tablenotes}
\end{threeparttable}
\end{table}

The 75--78\% reduction is the largest percentage reduction of any L-series metric.
Enacted map values of $\delta \approx +0.28$--$+0.40$ are dramatically above the
\textsc{bisect} baseline of $\delta \approx +0.07$--$+0.09$. If courts were to
develop interest in Declination, the gap of $\Delta\delta \approx 0.21$--$0.31$
provides strong evidence of deliberate cracking beyond what compactness-constrained
redistricting produces.

\subsection{NC District-Level Decomposition}

Table~\ref{tab:decl-nc-decomp} illustrates the geometric structure underlying
NC's Declination values by showing the average Democratic vote share in each
party's won districts under \textsc{bisect} and enacted plans.

\begin{table}[ht]
\centering
\caption{NC Declination Decomposition: Democratic and Republican Won-District Centroids}
\label{tab:decl-nc-decomp}
\begin{tabular}{@{}lrrrr@{}}
\toprule
Plan & $\bar{v}_D^+$ & $\bar{v}_R^-$ & $\theta_D$ & $\theta_R$ \\
\midrule
Bisect$^\dagger$ & $0.64$ & $0.44$ & $0.27$ & $0.21$ \\
Enacted 2022    & $0.71$ & $0.45$ & $0.66$ & $0.11$ \\
\bottomrule
\end{tabular}
\end{table}

The key difference is not in $\bar{v}_R^-$ (Republican-won district centroid), which
is similar for both plans ($0.44$--$0.45$), but in $\bar{v}_D^+$ (Democratic-won
district centroid), which is substantially higher for the enacted plan ($0.71$ vs.\
$0.64$). The enacted plan packs Democratic voters into fewer, higher-margin districts,
pulling $\bar{v}_D^+$ away from 50\% and increasing $\theta_D$. The cracking component
(low $\bar{v}_R^-$) is similar for both plans, reflecting that the geographic distribution
of Republican voters does not differ dramatically between compact and enacted boundaries.

This decomposition illustrates a key limitation of Declination: it does not distinguish
between packing (high $\bar{v}_D^+$) and cracking (low $\bar{v}_R^-$) as the source
of positive $\delta$. Both produce the same positive Declination via different mechanisms.
The NC evidence shows that the enacted map's elevated Declination is driven primarily
by packing (high $\bar{v}_D^+$) rather than cracking of Republican margins (the
Republican-won district centroid is similar in both plans).

\subsection{Election-Cycle Sensitivity}

Unlike MM (which is exactly invariant to uniform swings) and unlike EG (which varies
substantially with swings), Declination has moderate sensitivity to partisan swings
because uniform swings can move districts across the 50\% threshold, changing
$k_D$ and $k_R$ and potentially changing $\bar{v}_D^+$ and $\bar{v}_R^-$ discontinuously.

For NC, a 5-point uniform Democratic swing reduces enacted $\delta$ from $+0.35$
to $+0.28$ (because some Republicans-won districts at 47--50\% flip to Democratic,
changing the partition). \textsc{Bisect} $\delta$ falls from $+0.08$ to $+0.07$
(one district near 50\% flips). The qualitative gap persists under any realistic
swing scenario.
